\section{Conclusion}
\label{sec:conclusion}

In this project, we propose a novel method named heuristic conservative soft actor-critic (HCSAC) to solve the offline multi-task reinforcement learning problem. The method is developed based on the classic algorithms of conservative Q-learning (CQL) and soft actor-critic (SAC) to suppress distribution shifts in the offline setting. We further propose task embedding to allow the agent to learn shared representations across multiple tasks. In addition, we follow conservative data sharing (CDS) to propose heuristic relabeling to augment the offline dataset, so that the agent can benefit from the shared transitions while avoiding performance degradation. We conduct a series of experiments on the MuJoCo benchmark to verify the effectiveness of our method. The results show that our proposed methods in general bring superior performance and improve the robustness of the agent.

We should point out that it is not sufficient to evaluate our methods only on the MuJoCo benchmark. On one hand, the offline transitions are only shared across two tasks, where \texttt{walk} and \texttt{run} are both low-level locomotion skills, so it is unnecessary to share any part of the policy. In real-world scenarios, we actually expect the agent can learn from a large number of tasks to acquire those essential skills. On the other hand, MuJoCo is an environment with low-dimensional state spaces where the vectorized observations already contain all the information needed for robotic control, so the multi-task learning methods cannot be fully exploited. We assert that our methods will show more advantages in the high-dimensional vision-based environments, where the agent can utilize transitions from different tasks to learn better shared representations.